\clearpage\fancyhead[L]{\small PALLAS.}\hypertarget{native-325-lower}{}\label{native-325-lower}
\begin{center}\textbf{PALLAS.}\par\vspace{4pt}\hrule\vspace{4pt}\fontsize{10}{12}\selectfont Monatliche Correspondenz. Band XXII. S. 591\,..\,595. 1810 Dezember.\par\vspace{4pt}\hrule\vspace{4pt}\end{center}

\par\indent „Ich beeile mich, Ihnen, mein teuerster Freund, eine wichtige Nachricht mitzuteilen“, schrieb Professor \textsc{Gauss}. „Ich habe Ihnen neulich schon geschrieben, dass ich die Störungen berechne, die Pallas in den Jahren 1803\,..\,1811 durch Jupiter erfahren hatte. Wie schlecht sich alle sechs bisher beobachteten Oppositionen durch rein elliptische Elemente darstellen lassen, habe ich in der Abhandlung gezeigt, die ich Ihnen gegenüber schon einige Male erwähnt habe. Als ich die Arbeit über diese Störungen vor einigen Tagen abgeschlossen hatte, war es für mich deshalb außerordentlich befriedigend festzustellen, dass sich unter Berücksichtigung der Störungen alle sechs Oppositionen genauer in Einklang bringen ließen, als ich erwartet hatte. Hier ist die Tabelle der Abweichungen:

\begin{center}\fontsize{10}{12}\selectfont\setlength{\tabcolsep}{4pt}\renewcommand{\arraystretch}{1.15}\begin{tabular}{r|r|r}
 & \shortstack{Mittlere\\Länge} & \shortstack{Heliozentr.\\Breite}\\\hline
1803 & $+1''3$ & $-1''0$\\
1804 & $-3.8$ & $+4.1$\\
1805 & $+3.9$ & $+6.0$\\
1807 & $-3.3$ & $+3.9$\\
1808 & $+3.2$ & $-17.0$\\
1809 & $-1.4$ & $-3.9$\\
\end{tabular}\end{center}
\noindent (Die Breite für 1808 war bekanntlich sehr schlecht bestimmt.)

\SourceLeaf{326}\par\indent „Da sich die Elemente 1809 bei der heliozentrischen Konjunktion mit Jupiter \emph{sehr stark} geändert hatten, wollte ich wissen, wie weit meine Ephemeride im Oktoberheft der \emph{Monatl. Corresp.} wohl abweichen könnte. Ich berechnete daher aus meinen neuen Ergebnissen unter Berücksichtigung der Störungen eine Position. Dabei fand ich für den 9. Januar 1811 mittags in Göttingen

{\fontsize{10}{12}\selectfont\[\begin{array}{llr|l}\text{AE }\PallasSymbol & 151^\circ\;37' & \text{also }1^\circ\;40'\text{ weniger} & \raisebox{-0.6\baselineskip}{als in der Ephemeride.}\\\text{Decl.} & 22^\circ\;10'\text{ südl.} & \text{—}\quad 11'\text{ südl.} & \end{array}\]}

\noindent Ich habe gerade noch eine zweite Position berechnet, für den 26. Oktober 1810 mittags in Göttingen

{\fontsize{10}{12}\selectfont\[\begin{array}{llr|l}\text{AE }\PallasSymbol & 134^\circ\;55' & \text{also }45'\text{ weniger} & \raisebox{-0.6\baselineskip}{als die Ephemeride}\\\text{Decl.} & 14^\circ\;28'\text{ südl.} & \text{—}\quad 6'\text{ südl.} & \end{array}\]}

\Needspace{15\baselineskip}\noindent Wenn Sie selbst einzelne Positionen berechnen möchten, können Sie dazu die folgenden Elemente verwenden:

\begin{center}\fontsize{10}{12}\selectfont\setlength{\tabcolsep}{4pt}\renewcommand{\arraystretch}{1.15}\begin{tabular}{llr}
Epoche 1811 Göttingen & .\;.\;.\;.\;.\;. & $126^\circ\;32'\;52''$\\
tägliche mittlere trop. Bewegung & .\;. & $769''012$\\
Logarithmus der halben Achse & .\;.\;.\;. & $0.44277$\\
Exzentrizität & .\;.\;.\;.\;.\;.\;.\;.\;.\;. & $0.24162\ (=\operatorname{Sin}13^\circ\;18'\;56'')$\\
Perihel für die obige Epoche & .\;.\;.\;.\;. & $120^\circ\;55'\;5''$\\
Knoten & .\;.\;.\;.\;.\;.\;.\;.\;.\;.\;.\;. & $172\;32\;30$\\
Neigung & .\;.\;.\;.\;.\;.\;.\;.\;.\;.\;. & $34\;35\;55$\\
\end{tabular}\end{center}

\noindent Vor dem Abschicken kann ich noch eine Position hinzufügen, für den 15. Dez. 1810 mittags in Göttingen.

{\fontsize{10}{12}\selectfont\[\begin{array}{llr|l}\text{AE }\PallasSymbol & 149^\circ\;37' & \text{also }1^\circ\;15'\text{ weniger} & \raisebox{-0.6\baselineskip}{als die Ephemeride}\\\text{Decl.} & 21''\;37'\text{ südl.} & \text{—}\quad 12'\text{ südlicher} & \end{array}\]}

\par\indent „Ich bitte Sie nun, in der ersten klaren Nacht nach Pallas zu suchen. Hier war es bisher immer unmöglich. Sie können sich leicht vorstellen, wie ungeduldig ich die Bestätigung dieses bemerkenswerten Ergebnisses erwarte, das bei den neuen Planeten bisher noch nicht vorgekommen ist.

\par\indent „Durch Interpolation ergibt sich für die Ephemeride im Oktoberheft die folgende Korrektur:
\Needspace{12\baselineskip}
\begin{center}\fontsize{10}{12}\selectfont\setlength{\tabcolsep}{4pt}\renewcommand{\arraystretch}{1.15}\begin{tabular}{ll|r|l}
 &  & AE. & Decl.\\\hline
1810 den 26. Oktober & Mittag & $-45'$ & $6'$ südl.\\
--- den 7. November & Mitternacht & 49 & 8\\
--- den 20. November & Mittag & 56 & 10\\
--- den 2. Dezember & Mitternacht & 65 & 11\\
--- den 15. Dezember & Mittag & 75 & 12\\
--- den 27. Dezember & Mitternacht & 87 & 12\\
1811 den 9. Januar & Mittag & 101 & 11\\
\end{tabular}\end{center}

\par\indent So bemerkenswerte Ergebnisse und der Wunsch unseres hochverehrten Freundes waren Ansporn genug, mit größtmöglicher Sorgfalt nach Pallas zu suchen. Leider machte der seit sechs Wochen äußerst ungünstige Himmel fast alle unsere Bemühungen zunichte. Am 18. Dezember klarte es etwas auf, und um zwei Uhr morgens sahen wir im Gesichtsfeld des Fernrohrs am parallaktischen Instrument drei Sterne. Einer davon war höchstwahrscheinlich Pallas, doch eine eindeutige Beobachtung war unmöglich, da sich der Himmel bald wieder zuzog. Auch konnten wir weder bei \emph{Piazzi} noch in der \emph{Histoire céleste} eine Positionsbestimmung für einen dieser drei Sterne finden. Am 19. De\SourceLeaf{327}zember gelang uns endlich eine Meridianbeobachtung. Genau an der Stelle, an der die korrigierte Ephemeride Pallas angab, beobachteten wir einen Stern 7...8. Größe, der höchstwahrscheinlich Pallas war.

\begin{center}\fontsize{10}{12}\selectfont\setlength{\tabcolsep}{4pt}\renewcommand{\arraystretch}{1.15}\begin{tabular}{c|c|c|c}
1809 & \shortstack{mittlere\\Seeberger Zeit} & AR. app. $\PallasSymbol$ & \shortstack{Deklination\\austr.}\\\hline
Dezember 19 & $16^{\text{u}}\;9^{\text{m}}\;7^{\text{s}}3$ & $150^\circ\;20'\;13''5$ & $22^\circ\;58''7$\\
\end{tabular}\end{center}

\par\indent Die Deklination ließ sich freilich nur am Passageinstrument beobachten und ist auf die Minute unsicher. Die Rektaszension dagegen ist genau und stimmt mit der korrigierten Bestimmung ausgezeichnet überein. Im Fernrohr des Quadranten war Pallas überhaupt nicht sichtbar. Ein kleiner Stern in der Nähe von Pallas, für den wir die Position
\[\text{AR. }149''\;28'\;36''.\quad\text{Decl. austr. }21^\circ\;59'\]
\noindent fanden, war unverändert an seiner Stelle geblieben. Wir zögerten, Professor \textsc{Gauss} diese einzelne Beobachtung mitzuteilen, weil wir hofften, bald durch eine zweite bestätigen zu können, dass es wirklich Pallas war. Leider hat der Himmel uns diese Freude bisher nicht gestattet. Deshalb teilen wir ihm nun unsere vermutete Position von Pallas mit. \textsc{Gauss} teilt unsere Vermutung und schrieb uns darauf Folgendes:

\par\indent „Ich freue mich sehr, dass sich meine Rechnung über die Störungen von Pallas so schön ausgezahlt hat. Das Wetter hat es mir zwar bisher unmöglich gemacht, den Planeten aufzusuchen. Ich zweifle aber nicht mehr daran, dass der von Ihnen beobachtete Stern Pallas war. \textsc{Harding} hat an dieser Stelle keinen Stern auf seiner Karte, aber den anderen, den Sie am 18. Dez. ebenfalls gesehen haben, hat er eingetragen, obwohl er in der \emph{Histoire céleste} nicht vorkommt.“
